% \ifnotes
%     Recall that:
%     \begin{equation*}
%     a = (c_0, c_1, \dots c_{N-1}) \pmod{Q} =
%     \begin{pmatrix}
%         c_0 \bmod{q_0} & c_1 \bmod{q_0} & \dots & c_{N-1}\bmod{q_0} \\
%         c_0 \bmod{q_1} & c_1 \bmod{q_1} & \dots & c_{N-1}\bmod{q_1} \\
%         \vdots & \vdots & \ddots & \vdots \\
%         c_0 \bmod{q_k} & c_1 \bmod{q_k} & \dots & c_{N-1}\bmod{q_k}
%     \end{pmatrix}
%     \end{equation*}
% \fi

\subsection{Brakerski-Vaikuntanathan}
Although we cannot use \textit{digit} decomposition, RNS is itself based on a different decomposition: \textit{CRT} decomposition. This means we can decompose an RNS polynomial using the RNS primes themselves as a basis as shown in \cite{cryptoeprint:2016/510, cryptoeprint:2021/204}. % At level $1$ (s.t. $Q = Q_i$ in the notation of \cite{cryptoeprint:2021/204}) this becomes:
\begin{align}
    \label{eq:bv:d}
    D_{Q}({a}) &= \left(\imod{a\left(\frac{Q}{q_0}\right)^{-1}}{0},\dots,\imod{a\left(\frac{Q}{q_i}\right)^{-1}}{i}\right) \\
    \label{eq:bv:p}
    P_{Q}({a}) &= \left(\qmod{a\frac{Q}{q_0}},\dots,\qmod{a\frac{Q}{q_i}}\right)
\end{align}

\begin{lemma}
    \label{lemma:bv_gadget}
    The vectors \eqref{eq:bv:d} and \eqref{eq:bv:p} satisfy the gadget property \textcolor{red}{\autocite{cryptoeprint:2021/204}}.
    \begin{equation}
        \langle D_Q(a), P_Q(b)\rangle = a\cdot b \pmod{Q}
    \end{equation}
\end{lemma}

% \noindent\textbf{HPS Optimization}
It was later shown by Halevi, Polyakov and Shoup \cite{cryptoeprint:2018/117} that the factors can be moved into $P_Q(a)$ without violating Lemma \ref{lemma:bv_gadget}. Since $[(Q/q_i)^{-1}]_{q_i}[Q/q_i] = 1$ in tower $j = i$ and 0 in towers $j\not= i$, this optimization effectively makes the decomposition free as both $\hps{D}(a) \cong a_\mathsf{limbs}$ and $\hps{P}(a) \cong a_{\mathsf{limbs}}\cdot I_k$ contains the same information as $a$, and is only structured differently. For convenience we let $\hat{q}_i = {Q}/{q_i}$ so we can write the elements of $\hps{P}(a)$ as $a[\hat{q}_i^{-1}]_{q_i} [\hat{q}_i]_{Q}$.

\begin{align}
    \hps{D}(a) &= \left( \imod{a}{0}, \imod{a}{1}, \cdots, \imod{a}{i}\right)\\
    \hps{P}(a) &= \left( a\cdot\imod{\hat{q}_0^{-1}}{0} \cdot \qmod{\hat{q}_0}, \dots, a\cdot\imod{\hat{q}_\ell^{-1}}{i} \cdot \qmod{\hat{q}_\ell} \right)
\end{align}

\begin{itemize}
    \item This saves many unnecessary calculations in practice (see \autoref{sec:impl})
    \item The factors $[\hat{q}_i^{-1}]_{q_i} [\hat{q}_i]_{Q}$ are static and can be pre-calculated
\end{itemize}

% Maybe promote to sub/section after discussing all the main gadget types, because it can be applied to every gadget
\subsection{Double Decomposition}
\begin{itemize}
    \item By \autoref{the:gadget}, we can use these gadget vectors to define an external product. 
    \item The noise added scales with $\sum_{i}q_i \leq k\cdot \max_i\{q_i\}$ (\textcolor{red}{Cite error bounds from TFHE paper}) 
    \item Fine for keyswitching, but too large for some applications like internal products \eqref{def:internal_product}.
    \item To combat this, we can add a second layer of decomposition to each decomposed element using the classic gadget technique with $g = (1/\beta, \dots,1/\beta^{\ell-1})$.
    \item As each element is a ring element in $R_{q}$ for some small $q$ that fits inside a computer register/word, applying this method does not incur the performance penalty of using multi-precision integers.
    \item This can be represented as $\bv{D}(a) = \bv{P}(a\cdot g^{-1}(a))$ and $\bv{D}(a) = \bv{P}(a\cdot g)$ where $g = ($ and $g^{-1}(a)$ satisfies $g^{-1}(a)\cdot g = x$), but this results in a matrix $\bv{P}\in R^{2\times k}_q$. To fit this within our established framework, we simply flatten $R^{\ell\times k}_q \to R^{1\times k\ell}_q$.
\end{itemize}

\begin{corollary}[Correctness of Double Decomposition]
    Applying digit decomposition to each element of $\hps{P}(a)$ individually yields a valid external product by \autoref{the:gadget}. The proof changes only structurally when $P_{BV}(m)$ and $D_{BV}(n)$ are flattened.

    \begin{proof}[Proof Sketch]
        Let $\bv{P}(a) = \hps{P}(a\cdot g)$ and $D_{BV}(a) = D_{HPS}(a\cdot g^{-1})$ where $g = (1, \beta, \dots, \beta^{\ell-1})$ and $g^{-1}$ is the digit extraction of $n$ such that $g^{-1}(n)\cdot g = n$. Using the fact that $P_{HPS}(\cdot), D_{HPS}(\cdot)$ satisfies the exact gadget property \eqref{property:gadget}, we see that the same is true for $P_{BV}(\cdot), D_{BV}(\cdot)$:
        \begin{gather}
            \langle P_{BV}(m), D_{BV}(n)\rangle = \langle P_{HPS}(m\cdot g), D_{HPS}(n\cdot g^{-1})\rangle \pmod{Q}\\
            = (m\cdot g) \cdot (n\cdot g^{-1}) = m\cdot n \cdot (g^{-1}\cdot g) \pmod{Q} \\
            = m\cdot n \pmod{Q}
        \end{gather}

        $\therefore\quad\langle P_{BV}(m), D_{BV}(n)\rangle = m\cdot n \pmod{Q}$ 
    \end{proof}
\end{corollary}

\begin{remark}
    \textcolor{red}{Now $\ell$ covers $\max_i\{q_i\}$ and the RGSW has $k\cdot\ell$ rows. This is the same length required if $\ell$ covers the entire $Q$ as we would have $k$ less numbers but each number would be $k$ digits longer. For example, with $\ell_{small} = 2$ and $Q=q_0q_1 \implies k = 2$ we get $2\cdot k \cdot \ell_{small} =2\cdot 2\cdot 2 = 8$ RGSW rows and the noise is equivalent to having $\ell_{big} = 4 \implies 2\cdot\ell_{big} = 8$ RGSW rows over $Q$. The results is the same, the number of rows is the same, the noise is the same, its just faster to calculate over each $q_i$ individually.}
\end{remark}

Double decomposition 

\ifnotes
    \newpage
    \subsubsection{Notes}
    \noindent\textcolor{red}{\textbf{Assuming I prove \autoref{the:gadget}, the following is superfluous and could be moved to a different section (about RGSW in RNS specifically, or implementation details.)}}
    
    % \begin{remark}
    Now consider the $j$-th element of $P_{Q}(a)$ and recall that $Q = q_0q_1\dots q_k$. The factor $\hat{Q}_j = Q/q_j$ is a product of every $q_i$ except $q_j \; \therefore\;$ the $i$-th tower of $P_{Q}(a)[j]$ is a multiple of $q_j$ and becomes $0$ when $i\not= j$. Further, when $i = j$ then the factor $\hat{Q}_i^{-1}\cdot\hat{Q}_i = 1$. As a consequence, the $i$-th element of $P_Q(a)$ is a CRT polynomial where all towers are 0, except for the $i$-th tower which is equal to $a$'s $i$-th tower. 
     
    If we expand $P_Q(a)$ so that the elements are displayed in coefficient form \eqref{form:coef}, the entries along the diagonal of the matrix-form $P_Q(a)$ are exactly the towers of $a$, while the non-diagonal elements/towers are 0. 
    
    \begin{gather*}
        a_j\cdot \hat{Q}_i^{-1}\cdot\hat{Q}_i = \begin{cases}
            a & i = j \\
            0 & i \not= j
        \end{cases}
    \\
    \\
        \implies\quad P_Q(a) \cong 
        \left(
            \begin{array}{c|c|c|c} 
                [a]_{q_0} & 0 & \cdots & 0 \\ 
                0 & [a]_{q_1} & \cdots & 0 \\
                \vdots & \vdots & \ddots & \vdots \\
                0 & 0 & \cdots & [a]_{q_k} \\
            \end{array}
        \right)
    \\
        = (a\cdot e_0, a\cdot e_1, \dots, a\cdot e_{k-1})
        = a\cdot I_k^{\mathsf{RNS}} = a\cdot g 
    \\ 
    \\
        \text{where } [a]_{q_i} = ([c_0]_{q_0}, [c_1]_{q_0}, \dots, [c_{N-1}]_{q_i}) 
\end{gather*}
% \end{remark}


\begin{remark}
    \label{remark:hps}
    The 0-towers that are not on the diagonal of $P_Q(a)$ adds no information. Therefore, evaluating $P_Q(a)$ is free as it contains exactly the same information as $a$. The same is true for $D_Q(a)$ by definition. More on this in \autoref{sec:impl:hps}.
\end{remark}

\subsubsection{HPS External Product}
Let $e_i$ be the RNS basis polynomials, i.e. $e_0 = 1 \overset{coef}{\to} (1, 0, 0, \dots )$ and $e_1 = x \overset{coef}{\to} (0, 1, 0, \dots )$ and so on. We define the gadget vector $g = (e_0, \dots, e_1)$ as the RNS basis. Then the RGSW/external product definitions follow naturally, and we can chose $G^{-1}(x) = G^{-1}(x) = [D_Q(x_0) \;|\; D_Q(x_1)]$ as it satisfies $G^{-1}(x)\cdot G = x$. \textcolor{red}{The hidden link to HPS is that $P_Q(m) = m\cdot g$ aka. $C = Z + I_2\otimes P_Q(a)$ and $G^{-1}(x)\cdot G = x$ via CRT.}

%\ifnotes
    \begin{gather*}
        G = I_2\otimes e = \left(
        \begin{array}{c|c}
             e_0 & 0  \\
             \vdots & \vdots \\
             e_{k-1} & 0 \\
             \hline
             0 & e_0 \\
             \vdots & \vdots \\
             0 & e_{k-1} \\
        \end{array}
        \right), \quad
        G^{-1}(x) = 
        \begin{bmatrix}
            D_Q(x_0) \\ D_Q(x_1)
        \end{bmatrix}^\intercal
        % = 
        % \left(
        % \begin{array}{ccc|ccc}
        %     [x_0]_{q_0} & \cdots & [x_0]_{q_{k-1}} & [x_1]_{q_0} & \cdots & [x_1]_{q_{k-1}} \\
        % \end{array}
        % \right)
    \\
        G^{-1}(x)\cdot G = 
        \left(
        \begin{array}{ccc|ccc}
            [x_0]_{q_0} & \cdots & [x_0]_{q_{k-1}} & [x_1]_{q_0} & \cdots & [x_1]_{q_{k-1}}
        \end{array}
        \right) \cdot \left(
        \begin{array}{c|c}
             e_0 & 0  \\
             \vdots & \vdots \\
             e_{k-1} & 0 \\
             \hline
             0 & e_0 \\
             \vdots & \vdots \\
             0 & e_{k-1} \\
        \end{array}
        \right)
    \\
        = \left(
        \begin{array}{c|c}
            \sum\limits_{i=0}^{k-1} [x_0]_{q_i}\cdot e_i& \sum\limits_{i=0}^{k-1} [x_1]_{q_i}\cdot e_i
        \end{array}
        \right)
        \overset{CRT}{=} (x_0, x_1) = x
    \\
    \\
        \text{So the definitions \eqref{def:rgsw} and \eqref{def:external_product} are valid for this gadget.} \\
        \text{And the external product is } x\boxdot Y = G^{-1}(x)\cdot Y \pmod{Q} 
    \\ 
    \\
        G^{-1}(x) \cdot (Z + m_y\cdot G)                  \\ %\pmod{Q} \\
        G^{-1}(x)\cdot Z + m_y \cdot (G^{-1}(x)\cdot G)   \\ %\pmod{Q} \\
        \underset{\text{noisy }\rlwe(0)}{\underbrace{G^{-1}(x)\cdot Z}} \quad +\, \underset{\rlwe(m_y\cdot m_x)}{\underbrace{m_y \cdot x}} %\pmod{Q}
    \end{gather*}
% \fi

\textcolor{red}{\textbf{Basically this whole section boils down to "do BV on the towers instead of the full polynomial"..?}}

\subsubsection{Noise}
Scales by $k\cdot\max_i\|q_i\|/2$ instead of $Q/2$ ($\sum_i q_i = k\cdot2^{60}$ instead of $\prod_i q_i = 2^{60\cdot k}$)?

%\ifnotes
    \noindent The noise is coming from the $\rlwe(0)$ on the lhs and $m_y\cdot \epsilon$ on the lhs. The noise for the rhs is obviously 0 or $\epsilon$ since we assume the message $m_y\in \{0,1\}$ is binary. The lhs is bounded in the worst case by...
    
    \begin{gather*}
        G^{-1}(x)\cdot Z = 
        % \left(
        % \begin{array}{ccc|ccc}
        %     [x_0]_{q_0} & \cdots & [x_0]_{q_{k-1}} & [x_1]_{q_0} & \cdots & [x_1]_{q_{k-1}} \\
        % \end{array}
        % \right)
        \begin{pmatrix}
            [x_0]_{q_0} \\ 
            \vdots \\ 
            [x_0]_{q_{k-1}} \\
            [x_1]_{q_0} \\ 
            \vdots \\
            [x_1]_{q_{k-1}}
        \end{pmatrix}^\intercal
        \cdot
        \begin{pmatrix}
            a_0^z\cdot s + \epsilon_0 & a_0^z \\
            \vdots \\
            a_{2k-1}^z\cdot s + \epsilon_{2k-1} & a_{2k-1}^z
        \end{pmatrix}
    \\
    \\
        \leq 
        \begin{pmatrix}
            [x_0]_{q_0} \\ 
            \vdots \\ 
            [x_0]_{q_{k-1}} \\
            [x_1]_{q_0} \\ 
            \vdots \\
            [x_1]_{q_{k-1}}
        \end{pmatrix}^\intercal
        \cdot
        \begin{pmatrix}
            a_0\cdot s + \epsilon & a_0 \\
            \vdots \\
            a_{2k-1}\cdot s + \epsilon & a_{2k-1}
        \end{pmatrix}
        \text{ where } \epsilon = \max_i\|\epsilon_i\|
    \\
    \\
        \left(
        \begin{array}{cc}
             \sum\limits_{i=0}^{k-1} [x_0]_{q_0} &  \\
             & 
        \end{array}
        \right)
    \end{gather*}
    
    % \newpage
    
    % Actual calculation:
    
    % \begin{gather*}
    %     G^{-1}(x) \cdot Y 
    % \\
    % \\
    %     (
    %     \begin{array}{ccc|ccc}
    %         [x_0]_{q_0} & \cdots & [x_0]_{q_{k-1}} & [x_1]_{q_0} & \cdots & [x_1]_{q_{k-1}}
    %     \end{array}
    %     )
    %     \left(
    %     \begin{array}{c|c}
    %         a_0\cdot s + m^{(0)} + \epsilon_0 & a_0 \\
    %         \vdots & \vdots \\
    %         a_{k-1}\cdot s + m^{(k-1)} + \epsilon_{k-1} & a_{k-1} \\
    %         \hline
    %         a_k\cdot s + \epsilon_k & a_k + m^{(0)} \\
    %         \vdots & \vdots \\
    %         a_{2k-2}\cdot s + \epsilon_{2k-2} & a_{2k-2} + m^{(k-1)} \\
    %     \end{array}
    %     \right)
    % \\
    % \\
    %     = (\sum_{i=0}^{k-1} [x_0]_{q_i} \cdot (z_0 + [m]_{q_i})
    % \end{gather*}
    
    % \newpage
%\fi
    
    \subsubsection{Second Layer of Decomposition II}
    The noise added by this gadget is still too large. The BV decomp. decreases the noise to $B/2$ and can be applied to each of the RNS-digits. I.e. instead of decomposing the full polynomial $a$ into digits in base $B$, we decompose each limb/tower. The gadget decomposition parameter $\ell = \lfloor\log(\max_i\{q_i\})/\log(B)\rfloor + 1$ is now such that the $B$-digits cover just $q_i$ instead of the full $Q$, but we need $\ell$ digits \textit{per limb}, so this is really equivalent to BV in a non-RNS setting!
    
    %%%
%%% Powers of base
%%%
\begin{algorithm}
\caption{Powers of Base}
\label{alg:powers-of-base}
\begin{algorithmic}[1]
\Require Evaluation form $a \in R_Q$, base $B = 2^b$, decomposition parameter $\ell$
\Ensure Vector of $a$ scaled $\hat{a} = (a, a B, \dots, aB^{\ell-1})$

\State Initialize an array of polynomials $\hat{a}$ of size $\ell$
\State Initialize a counter $B' \leftarrow 1$ \Comment{Calculate $B^i\bmod{q_j}$ offline}

\For{$i = 0$ to $\ell - 1$} % \textbf{in parallel}
    \State $\hat{a}_i \leftarrow a \cdot B' \pmod{Q}$ \Comment{Scale polynomial}
    \State $B' \leftarrow B' \cdot B \pmod{Q}$ \Comment{Update counter}
\EndFor

\State \Return $\hat{a}$
\end{algorithmic}
\end{algorithm}
    
    Let $\hat{m} = (m, mB, \dots, mB^{\ell-1})$ be generated by \autoref{alg:powers-of-base}. The gadget vector becomes
    
    % Then the RGSW ciphertext is the familiar TFHE form, except that we consider $Z$ to be of size $2\cdot k\cdot \ell$ where there are $k$ moduli $q_0q_1\cdots q_{\ell-1}$, each one is expressable by $\ell$ digits in base $B$:
    % \begin{equation*}
    %     C = Z + m\cdot G = Z + P_Q(\hat{m})\otimes I_2
    % \end{equation*}
    
    \newpage
\fi

\ifnotes
    \newpage
    \newpage
    \newcommand{\coef}[2]{c_{#1}\bmod{q_{#2}}}
    
    \begin{gather*}
        \text{Recall that } \mathbf{a} = 
        \begin{pmatrix}
            \mathbf{a}\bmod{q_0} \\
            \mathbf{a}\bmod{q_1} \\
            \vdots \\
            \mathbf{a}\bmod{q_{k-1}}
        \end{pmatrix}
        = 
        \begin{pmatrix}
            \coef{0}{0} & \coef{1}{0} & \dots & \coef{N-1}{0} \\
            \coef{0}{1} & \coef{1}{1} & \dots & \coef{N-1}{1} \\
            \vdots & \vdots & \ddots & \vdots \\
            \coef{0}{k-1} & \coef{1}{k-1} & \dots & \coef{N-1}{k-1}
        \end{pmatrix} 
    \\
    \\
        \text{Examine the $j$-th element of } P_Q(\mathbf{a}) \to a\cdot\imod{\left(\frac{Q}{q_j}\right)^{-1}}{j} \cdot \qmod{\frac{Q}{q_j}}
    \\ 
    \\
        \text{Let } \hat{Q}_j = \frac{Q}{q_j}, \text{ so that } P_Q(\mathbf{a})_{(j=0)} = 
        \begin{pmatrix}
            \left[\mathbf{a}\cdot \hat{Q}_0^{-1}\cdot \hat{Q}_0\right]_{q_0} \\
            \left[\mathbf{a}\cdot \hat{Q}_0^{-1}\cdot \hat{Q}_0\right]_{q_1} \\
            \vdots \\
            \left[\mathbf{a}\cdot \hat{Q}_0^{-1}\cdot \hat{Q}_0\right]_{q_{k-1}}
        \end{pmatrix}
    \\
    \\
        \text{When the tower index $i$ is } 
        \begin{cases}
            i \not= j,      & \hat{Q}_j \equiv Q/q_j \equiv  ({q_0\cdots q_{i}\cdots q_{k-1}})/q_j \equiv n\cdot q_i \equiv 0 \pmod{q_i} \\
            i = j,  & \hat{Q}_j^{-1}\hat{Q}_j \equiv 1 \pmod{q_i} 
        \end{cases}
    \\
    \\
        \therefore P_Q(\mathbf{a}) = 
        \begin{pmatrix}
            [\mathbf{a}]_{q_0}  \\
            & [\mathbf{a}]_{q_0} \\
            & & \ddots \\
            & & & [\mathbf{a}]_{q_{k-1}}
        \end{pmatrix} = 
        \begin{pmatrix}
            \begin{pmatrix}
                \mathbf{a}\bmod{q_0} \\
                0 \\
                \vdots \\
                0
            \end{pmatrix} & \mathbf{0} & \mathbf{0} & \mathbf{0} \\
            \mathbf{0} & \begin{pmatrix}
                0 \\
                \mathbf{a}\bmod{q_1} \\
                \vdots \\
                0
            \end{pmatrix} & \mathbf{0} & \mathbf{0} \\
            \mathbf{0} & \mathbf{0} & \ddots & \vdots \\
            \mathbf{0} & \mathbf{0} & \cdots & \begin{pmatrix}
                0 \\
                0 \\
                \vdots \\
                \mathbf{a}\bmod{q_{k-1}}
            \end{pmatrix} \\
        \end{pmatrix}
    \end{gather*}
    
    \let\coef\undefined
    
    \newpage
    
    In other words, $P_Q(\mathbf{a})$ is a $(k-1)\times (k-1)$ sparse matrix where the $i$-th diagonal element is the $i$-th tower of $\mathbf{a}$. Let $\mathbf{a}^{(i)}$ denote the RNS polynomial where every tower is 0 except the $i$-th tower, so that $P_Q(a) = (\mathbf{a}^{(0)}, \mathbf{a}^{(1)}, \dots, \mathbf{a}^{(k-1)})^\intercal$. Let $g = $ Then its clear that the external product $Z + mG = Z + I_2\otimes P_Q(\mathbf{a})$ has $2k$ rows.
    
    \begin{gather*}
        \left(
        \begin{array}{c|c}
            a_0\cdot s + m^{(0)} + \epsilon_0 & a_0 \\
            \vdots & \vdots \\
            a_{k-1}\cdot s + m^{(k-1)} + \epsilon_{k-1} & a_{k-1} \\
            \hline
            a_k\cdot s + \epsilon_k & a_k + m^{(0)} \\
            \vdots & \vdots \\
            a_{2k-2}\cdot s + \epsilon_{2k-2} & a_{2k-2} + m^{(k-1)} \\
        \end{array}
        \right)
        =
        \begin{pmatrix}
            \rlwe_s(m^{(0)}) \\
            \vdots \\
            \rlwe_s(m^{(k-1)}) \\
            \rlwe_s(-s\cdot m^{(0)}) \\
            \vdots \\
            \rlwe_s(-s\cdot m^{(k-1)})
        \end{pmatrix}
    \end{gather*}
    
    Now take the decomposition $D_Q(\mathbf{x}) = ([\mathbf{x}]_{q_0}, [\mathbf{x}]_{q_1}, \dots, [\mathbf{x}]_{q_{k-1}})$. Again, the elements are simply the RNS polynomials with a single non-zero tower, $D_Q(\mathbf{x}) = (\mathbf{x}^{(0)}, \mathbf{x}^{(1)}, \dots, \mathbf{x}^{(k-1)})$. % And we have the external product using $G^{-1}(\mathbf{x}) = I_2 \cdot (D_Q(a\cdot s + m_x + \epsilon),\; D_Q(a))$.
    
    % \begin{equation}
    %     x\boxdot Y = G^{-1}(x)\cdot (Z + I_2\otimes P_Q(m_y)) \pmod{Q}
    % \end{equation}
    
    % Calculated:
    % \begin{gather*}
    %     D_Q(\mathbf{x}) = (\mathbf{x}^{(0)}, \mathbf{x}^{(1)}, \dots, \mathbf{x}^{(k-1)}) \\
    %     P_Q(\mathbf{y}) = (\mathbf{y}^{(0)}, \mathbf{y}^{(1)}, \dots, \mathbf{y}^{(k-1)}) \\
    % \\
    % \\
    %     \text{Let } \mathbf{x} = (\mathbf{x}_0, \mathbf{x}_1) = (a\cdot s + m_x + \epsilon, a) \\
    % \\
    % \\
    %     G^{-1}(\mathbf{x}) = 
    %     \begin{pmatrix}
    %         D_Q(\mathbf{x}_0 = a\cdot s + m_x + \epsilon) & 0 \\    
    %         0 & D_Q(\mathbf{x}_1)
    %     \end{pmatrix}
    %     =
    %     \left(
    %     \begin{array}{c|c}
    %          \mathbf{x}_0^{(0)} & 0  \\
    %          \vdots & \vdots \\
    %          \mathbf{x}_0^{(k)} & 0 \\
    %          \hline
    %          0 & \mathbf{x}_1^{(0)}  \\
    %          \vdots & \vdots \\
    %          0 & \mathbf{x}_1^{(k)} \\
    %     \end{array}
    %     \right)
    % \end{gather*}
    
    % \begin{gather*}
    %     x\boxdot Y =
    %     \left(
    %     \begin{array}{c|c}
    %          \mathbf{x}_0^{(0)} & 0  \\
    %          \vdots & \vdots \\
    %          \mathbf{x}_0^{(k)} & 0 \\
    %          \hline
    %          0 & \mathbf{x}_1^{(0)}  \\
    %          \vdots & \vdots \\
    %          0 & \mathbf{x}_1^{(k)} \\
    %     \end{array}
    %     \right) \cdot \left(
    %     \begin{array}{c|c}
    %         a_0\cdot s + m^{(0)} + \epsilon_0 & a_0 \\
    %         \vdots & \vdots \\
    %         a_{k-1}\cdot s + m^{(k-1)} + \epsilon_{k-1} & a_{k-1} \\
    %         \hline
    %         a_k\cdot s + \epsilon_k & a_k + m^{(0)} \\
    %         \vdots & \vdots \\
    %         a_{2k-2}\cdot s + \epsilon_{2k-2} & a_{2k-2} + m^{(k-1)} \\
    %     \end{array}
    %     \right) 
    % \\
    % \\
    %     \left(
    %     \begin{array}{c|c}
    %          (a\cdot s + m_x + \epsilon)^{(0)} & 0  \\
    %          \vdots & \vdots \\
    %          (a\cdot s + m_x + \epsilon)^{(0)} & 0 \\
    %          \hline
    %          0 & a^{(0)}  \\
    %          \vdots & \vdots \\
    %          0 & a^{(k)} \\
    %     \end{array}
    %     \right) \cdot \left(
    %     \begin{array}{c|c}
    %         a_0\cdot s + m^{(0)} + \epsilon_0 & a_0 \\
    %         \vdots & \vdots \\
    %         a_{k-1}\cdot s + m^{(k-1)} + \epsilon_{k-1} & a_{k-1} \\
    %         \hline
    %         a_k\cdot s + \epsilon_k & a_k + m^{(0)} \\
    %         \vdots & \vdots \\
    %         a_{2k-2}\cdot s + \epsilon_{2k-2} & a_{2k-2} + m^{(k-1)} \\
    %     \end{array}
    %     \right) 
    % \\
    % \\
    %     \text{Note that for any polynomial $p$ we have } \sum_i \mathbf{p}^{(i)} = \mathbf{p}
    % \\
    % \\
    %     = \left(
    %     \begin{array}{c|c}
    %         \sum_i \mathbf{x}^{(i)}_0 \cdot (a_i\cdot s + m_y^{(i)} + \epsilon_0) & 
    %         \dots \\
    %         \hline
    %         \dots & 
    %         \dots
    %     \end{array}
    %     \right)
    %     = \left(
    %     \begin{array}{c|c}
    %         (\sum_i \mathbf{x}^{(i)}_0) \cdot (\sum_i a_i\cdot s + m_y^{(i)} + \epsilon_i) & 
    %         \dots \\
    %         \hline
    %         \dots & 
    %         \dots
    %     \end{array}
    %     \right)
    % \\
    %     = \left(
    %     \begin{array}{c|c}
    %         \mathbf{x}_0 \cdot [(\sum_i a_i\cdot s + \epsilon_i) + m_y] & 
    %         \dots \\
    %         \hline
    %         \dots & 
    %         \dots
    %     \end{array} 
    %     \right)
    % % \\
    % % \\
    % %     = \left(
    % %     \begin{array}{c|c}
    % %         \sum_i \mathbf{x}^{(i)}_0 \cdot \mathbf{y}^{(i)}_0 & \sum_i \mathbf{x}^{(i)}_0 \cdot \mathbf{y}^{(i)}_0 \\
    % %         \hline
    % %         %\sum_i \mathbf{x}^{(i)}_0 \cdot \mathbf{y}^{(i)}_0 & \sum_i \mathbf{x}^{(i)}_0 \cdot \mathbf{y}^{(i)}_0
    % %     \end{array}
    % %     \right)
    % \end{gather*}
    
    \newpage
    
    %%%
    %%% Encrypt
    %%%
    \begin{algorithm}
    \caption{Encrypt RGSW}
    \label{alg:rgsw:encrypt0}
    \begin{algorithmic}[1]
    \Require Polynomial $m$ and param $\ell$
    \Ensure RGSW as matrix of $2\times k \ell$ polynomials
    
    \State Initialize a vector $\mathbf{C}$ of size $2 \times k \ell$
    
    \For{$l \in [0,1]$}
        \For{$j \in [0, k)$}
            \For{$i \in [0, \ell)$}
                \State Create fresh $z \leftarrow \mathsf{RLWE}(0) = (c_0, c_1)$
                \State $c_l \leftarrow c_l + (0, \dots, 0, [mB^i]_{q_j}, 0, \dots, 0)$
                \State $\mathbf{C}[l][k\cdot \ell + i] = z$
            \EndFor
        \EndFor
    \EndFor
    
    \State \Return $\mathbf{C}$
    \end{algorithmic}
    \end{algorithm}
    
    \begin{remark}
        Due to the sparsity of $P_Q$ and $k$ repetitions, we can optimize the storage and parallallize this function by unrolling the loop, leading to \autoref{alg:rgsw:encrypt}. % Note: (maybe) the modulo operations can be further optimized with bitwise operations assuming $k\ell = 2^x$ is a power of 2?
        This version only adds the non-zero towers. We can also skip generating the zero-towers in the first place, as described in \autoref{sec:impl:hps}.
    \end{remark}
    
    % \newpage
    The inverse of the power of base is the (signed) digit decomposition:
    
    %%%
%%% Signed digit decomposition (optimized for B = 2^b)
%%%
\begin{algorithm}
\caption{Signed LSB Base-$B$ Decomposition}
\label{alg:lsb-decomp}
\begin{algorithmic}[1]
\Require Coefficient $a \in \mathbb{Z}$, base $B = 2^b$, number of digits $\ell$
\Ensure Vector of signed digits $\mathbf{d} = (d_0, d_1, \dots, d_{\ell-1})$ where $d_j \in [-B/2, B/2)$
\State $a' \gets a$
\For{$j = 0$ to $\ell - 1$}
    \State $d_j \gets a' \pmod B$ \Comment{Extract the lower bits}
    \If{$d_j \ge B/2$}
        \State $d_j \gets d_j - B$ \Comment{Shift to signed range}
    \EndIf
    \State $a' \gets \frac{a' - d_j}{B}$ \Comment{Shift down and apply the carry}
\EndFor
\State \Return $\mathbf{d}$
\end{algorithmic}
\end{algorithm}

% function ExtractDigit(poly_coeff, j, b_bits):
%     mask = (1 << b_bits) - 1
%     shift = j * b_bits
%     result_poly = new Polynomial()
    
%     # Iterate through the N coefficients of the polynomial
%     for x from 0 to N - 1:
%         # Standard bitwise shift and mask to isolate the digit
%         result_poly[x] = (poly_coeff[x] >> shift) & mask
        
%     return result_poly
    
    Which is encoded into the external product:
    \begin{equation*}
        a\boxdot B = D_Q(a) \cdot B = D_Q(a)\cdot (Z + m_b\cdot G)
    \end{equation*}
    
    \noindent\textbf{Noise Growth (revisited)}
    Manageable but not optimal. It is now bounded by $B/2$, but performance is still poor as $\ell$ is quite high ($\ell$ digits must cover $\max_i{\|q_i\|}$).

\newpage
\newpage
\fi